\section{Related Work}
\label{sec:related}

\paragraph{Deep ensembles.}
Lakshminarayanan et al.\ \cite{lakshminarayanan2017simple} introduced deep
ensembles as a simple uncertainty-estimation and accuracy-boosting method:
train $N$ independent models with different random seeds and average their
softmax outputs. Deep ensembles are strong baselines for both accuracy and
calibration; they remain widely used despite their $N\times$ inference cost.
Our method directly extends deep ensembles by adding \emph{trajectory-length}
as a second diversity axis on top of seed.

\paragraph{Snapshot ensembles.}
Huang et al.\ \cite{huang2017snapshot} proposed collecting multiple
checkpoints from a single training run via cyclic learning rates. Each
cycle restarts the LR, allowing the model to explore different local
minima during one training. Our method differs: we use \emph{two separate}
training runs (not one with cycles) and vary only the cosine endpoint.
This is conceptually simpler and does not require a custom LR schedule.
Empirically, on CIFAR-10 with matched compute, our trick gives comparable
or larger gains than 5-snapshot ensembling (see \cref{sec:experiments}).

\paragraph{Stochastic Weight Averaging (SWA).}
Izmailov et al.\ \cite{izmailov2018averaging} proposed averaging weights
(not outputs) collected along a single training trajectory after the LR
has decayed. SWA has become a standard ``free'' trick for PyTorch
practitioners via \texttt{torch.optim.swa\_utils}. Our method differs in two
ways: (i)~we average \emph{outputs}, not \emph{weights} (better calibration,
preserves ensemble uncertainty); (ii)~we use \emph{two separate trainings},
which lets us vary diversity across seeds and trajectory lengths
simultaneously. Head-to-head comparison on CIFAR-10 shows that \ours{} is
competitive with SWA.

\paragraph{Fast Geometric Ensembling (FGE).}
Garipov et al.\ \cite{garipov2018loss} showed that different local minima
in a deep network's loss landscape can be connected by simple Bezier
curves, and proposed FGE: walk along these curves and collect snapshots.
FGE provides strong diverse ensembles at the cost of a complex training
procedure that requires a pretrained initial model plus a curve-finding
loss. \ours{} achieves comparable diversity by simply varying the training
endpoint, without any architecture or loss changes.

\paragraph{BatchEnsemble and hyperparameter ensembles.}
Wen et al.\ \cite{wen2020batchensemble} introduced efficient ensemble
parameterizations using rank-1 perturbations of a base model. Wenzel et
al.\ \cite{wenzel2020hyperparameter} explored training an ensemble where
each member varies a hyperparameter (dropout, weight decay, etc.). Our
work is complementary: trajectory length is a new axis of diversity that
composes with any of these methods.

\paragraph{EEG emotion classification.}
FACED \cite{liu2024faced} and SEED-V \cite{zheng2019investigating} are
two standard EEG emotion classification benchmarks. Recent work on EEG
foundation models---LaBraM \cite{jiang2024labram}, CBraMod
\cite{wang2025cbramod}, BIOT \cite{yang2023biot}, REVE
\cite{kostas2025reve}, EMOD \cite{chen2026emod}---has driven steady
improvements through large-scale pretraining. Our new SOTA on FACED
($\facedacc$) is achieved not by a new pretraining or architecture
contribution but by applying our general ensemble trick to the EMOD
backbone. We view our EEG result as an \emph{application demonstration}
of the ensemble trick rather than a domain-specific contribution.

\paragraph{EEG KD with LLM teachers.}
Several recent works use LLM-derived emotion vectors as KD teachers for
EEG models \cite{zhang2024emotionkd, li2025emod, aqa2024emotion}. Our
previous work showed that the LLM's \emph{content} is not the source of
the signal---a random 9-dimensional orthonormal tensor matches LLM-derived
ones within noise---and that the mechanism is ``fixed-prototype
regularization''. We use this finding in the EMOD d6 recipe but it is
orthogonal to the trajectory-length ensemble trick that is the focus of
this paper.
